\documentclass[10pt]{article}

\usepackage[utf8]{inputenc}

\usepackage[T1]{fontenc}

\usepackage{amsmath}

\usepackage{amsfonts}

\usepackage{amssymb}

\usepackage[version=4]{mhchem}

\usepackage{stmaryrd}

\usepackage{bbold}



\begin{document}

вание следует из монотонного возрастания и ограниченности сверху последовательности



\[

1+\frac{1}{2}+\ldots+\frac{1}{n}-\ln (n+1)

\]



Арифметич. природа Э. п. не изучена, неизвестно (1984) даже является ли она рациональным числом или нет. ЭИЛЕРА ПРЕОБРАЗОВАНИЕ -1) Л. Д. Кудрявчев. ЭЙЛЕРА ПРЕОБРАЗОВАНИЕ - 1) Э. П. р я д о в: если дан числовой ряд



\[

\sum_{n=0}^{\infty}(-1)^{n} a_{n}

\]



то ряд



\[

\sum_{n=0}^{\infty} \frac{\Delta^{n} a_{0}}{2^{n+1}}

\]



наз. рядом, полученным из ряда (1) Э. ш. рядов. Здесь



\[

\Delta^{n} a_{0}=\sum_{k=0}^{n}(-1)^{k} C_{n}^{k} a_{k}

\]



Если ряд (1) сходится, то сходится и ряд (2) и притом к той же сумме, что и ряд (1). Если ряд (2) сходится (в этом случае ряд (1) может расходиться), то ряд (1) наз. с у м и и р у е м м п о Ә й ле р у.



Если ряд (1) сходится, $a_{n}>0$, для всех $k=0,1,2, \ldots$ последовательность



\[

\Delta^{k} a_{n}=\sum_{l=0}^{k}(-1)^{l} C_{k}^{l} a_{n+l}

\]



монотонная и



\[

\frac{a_{n+1}}{a_{n}} \geqslant q>\frac{1}{2},

\]



то сходимость ряда (2) быстрее сходимости ряда (1) (см. Сходимость). Л. Д. Кудрявчев,



\begin{enumerate}

  \setcounter{enumi}{1}

  \item Э. п. - интегральное преобразование вида

\end{enumerate}



\[

w(z)=\int_{C}(z-t)^{\alpha} v(t) d t

\]



где $C$ - контур в комплексной плоскости $t$. Предложено Л. Эйлером (L. Euler, 1769).



Ә. ш. применяется к линейным обыкновенным дифференциальным уравнениям



\[

L w=\sum_{j=0}^{n}(-1)^{j} w^{(j)} \sum_{k=0}^{n-j} C_{n+\beta-j-1}^{n-k-j} Q_{j}^{(n-k-i)}(z)=0,

\]



где $Q_{j}(z)$ - многочлен степени $\leqslant n-j$ и $\beta-$ константа. В таком виде можно представить любое линейное уравнение



\[

P_{n}(z) w^{(n)}+P_{n-1}(z) w^{(n-1)}+\ldots+P_{0}(z) w=0,

\]



где $P_{j}(z)$ - многочлены степени $\leqslant j$, степень $P_{n}(z)$ равна $n$. У равнение



\[

M v \equiv \sum_{j=0}^{n}(-1)^{j}\left(Q_{n-j}(z) v\right)^{(j)}=0

\]



наз. П ре о б р а з в ани им Ә йле р а уравнения (2). Если $w(z)$ определена формулой (1), причем $\alpha=\beta+n-1$, то справедливо тождество



\[

L w=\int_{C}(z-t)^{\alpha} M(v) d t

\]



при условии, что внеинтегральная подстановка, к-рая возникает при интегрировании по частям, обращается в нуль. Отсюда видно, тто если $M(v)=0$, то $\boldsymbol{w}(z)-$ решение уравнения (2).



Э. п. позволяет понизить порядок уравнения (2), если $Q_{j}(z) \equiv 0$ при $j>q, q<n$. При $q=0, q=1$ уравнение (2) интегрируется (см. Похгаммера уравнение).



Лuт.: [1] А й н с Э. Л., Обыкновенные дифференциальные уравнения, пер. с англ., Харьков, 1939; [2] К а м к е Э., Справочник по обыкновенным диффференциальным уравнениям, пер. с нем., 5 изд., М., $1976 . \quad$ М. В. Федорюк. 3) Э. п. - 1 - г о р о д а - интегральное преобразование вида



\[

I \mu_{f}(x)=\mathfrak{N}_{\mu}\{f(t) ; x\}=\frac{1}{\Gamma(\mu)} \int_{0}^{x} f(t)(x-t) \mu-1 d t,

\]



где $\mu, x$ - комплексные переменные, причем путем ннтегрирования является отрезок $t=x \tau, 0<\tau<1$.



Ә. Л. 1-го рода наз. также дробным и н т е г р а л о м $\mathrm{P}$ и м а н а - Л и у в и л л я порядка $\mu$. (Иногда под интегралом Римана - Лиувилля понимают интеграл



\[

\frac{1}{\Gamma(\mu)} \int_{x}^{a} f(t)(t-x)^{\mu-1} d t=\mathfrak{R}_{\mu}\{f(a-t) ; a-x\},

\]



где $a$ - комплексное число.)



При нек-рых условиях на функции $f(x), g(x)$ имеют место следующие равенства:



\[

I^{0} f(x)=f(x)

\]



$I \mu[\alpha f(x)+\beta g(x)]=\alpha \boldsymbol{I}_{f}(x)+\beta I^{\mu} g(x)$, $\alpha, \beta$-комплексные постоянные,



\[

I \mu\left[I^{v} f(x)\right]=I^{\mu+v} f(x),

\]



$I^{n} f(x)=\int_{0}^{x} d t_{1} \int_{0}^{t_{1}} d t_{2} \ldots \int_{0}^{t_{n-2}} d t_{n-1} \int_{0}^{t_{n-1}} f\left(t_{n}\right) d t_{n}$,



\[

I-n_{f}(x)=\frac{d^{n}}{d x^{n}} f(x), \quad n=1,2,3, \ldots

\]



Э. П. 2-го р о д а - интегральное преобразованиө вида



$K^{\mu} f(x)=\mathbb{M}_{\mu}\{f(t) ; x\}=\frac{1}{\Gamma(\mu)} \int_{x}^{\infty} f(t)(t-x) \mu-1 d t$, где $\mu, x$ - комплексные переменные, причем путем интегрирования является луч $t=x \tau, \tau>1$, или $t=x+\tau$, $\tau>0$. При нек-рых условиях имеют место следующие равенства



\[

K \circ f(x)=f(x)

\]



\[

K^{\mu}[\alpha f(x)+\beta g(x)]=\alpha K^{\mu_{f}}(x)+\beta K^{\mu^{\prime}} g(x),

\]



$\alpha, \beta$-комплексные постоянные,



\[

K^{\mu}\left[K^{v} f(x)\right]=K^{\mu+v_{f}}(x),

\]



$K^{n} f(x)=\int_{x}^{\infty} d t_{1} \int_{t_{1}}^{\infty} d t_{2} \ldots \int_{t_{n-2}}^{\infty} d t_{n-1} \int_{t_{n-1}}^{\infty} f\left(t_{n}\right) d t_{n}$,



\[

K-n_{f}(x)=\frac{d^{n}}{d x^{n}} f(x), \quad n=1,2,3, \ldots

\]



Э. п. 2-го рода иногда наз. дробным и н те г р а л о м В е й л я порядка $\mu$.



Указанные преобразования введены также для обобщенных функций.



Лum.: 'Бр ы ч к о в Ю. А., П р у Д н и ко в А. П., Интегральные преобразования обобщенных функций, М., 1977.



ЭЙЛЕА ПРОИЗВЕДЕНЙЕ - бесконечное произвдение вида



\[

\prod_{p}\left(1-\frac{1}{p^{s}}\right)^{-1}

\]



где $s$ - действительное число и $p$ пробегает все простые числа. Это произведение абсолютно сходится при всех $s>1$. Аналогичное произведение для комплексных чисел $s=\sigma+i t$ абсолютно сходится при $\sigma>1$ и задает в этой области дзета-функцию Римана



\[

\zeta(s)=\prod_{p}\left(1-\frac{1}{p^{s}}\right)^{-1}=\sum_{n=1}^{\infty} \frac{1}{n^{s}} .

\]



С. А. Степанов.



ӘЙЛЕРА ПРЯМАЯ - прямая, на к-рой лежат точка $H$ пересечения высот треугольника, точка $S$ пересечения медиан и точка $O$ - центр описанной окружности. Если Э. п. проходит через вершину треугольника, то треугольник либо равнобедренный, либо прямоугольный, либо прямоугольный равнобедренный. Для отрезков Э. п. выполняется соотношение:



\[

O H: S H=1: 2

\]





\end{document}